% !TEX root = ../thesis.tex

\chapter{Discussion}
\label{ch:discussion}

% Where the Results chapter reports numbers, this chapter interprets them,
% situates the findings against prior work, and acknowledges limits.

\section{Interpretation of Results per Research Question}
\label{sec:discussion:interpretation}

\subsection{RQ1: Complementary Evaluation Metrics}

The clearest methodological finding from RQ1 is that a
distance-based privacy result cannot be interpreted as evidence of
protection when the synthetic data have poor fidelity. PATE-GAN is
a direct example for this case. It produced large nearest-neighbour distances and zero
outlier linkability. From this evidence alone one can conclude that PATE-GAN
is the obvious choice, but when you include fidelity analysis it shows substantial distributional
differences and weak predictive utility. Its distance results therefore
reflected poor similarity to the real data rather than a
favourable privacy and utility balance.

The other privacy diagnostics also had different roles. Membership
inference and NNAA privacy loss provided little separation between the
generators under the selected attacks. NNDR and attribute inference
identified some model-specific differences. These findings show why
privacy metrics should be interpreted together and alongside utility
results. This supports previous recommendations for multidimensional
synthetic-data evaluation
\cite{hernandez2022synthetic,lautrup2025syntheval}.

The conclusions also depend on how the metrics were implemented. Mean KS
was applied to nominal variables represented by numerical codes.
Correlation difference measured only pairwise Pearson relationships
between encoded columns. The privacy results depend on the selected
attacks, attributes, distance measure, and outlier definition. Future
work should use type-specific marginal distances, mixed-data association
measures, repeated runs, and stronger privacy attacks. These extensions
would show whether the current findings remain stable under different
evaluation choices.

\subsection{RQ2: Comparison of the Evaluated GAN and LLM Configurations}

The main finding from RQ2 is that utility separated the evaluated
configurations more clearly than the selected privacy attacks. The two
non-private Pythia configurations produced the highest observed
predictive-performance results. Pythia-1B also produced the lowest
pairwise linear-correlation difference. HealthGAN produced the lowest
observed mean KS and the highest Cox confidence-interval overlap.
Membership inference and outlier linkability provided less separation,
but NNDR and attribute inference identified some model-specific
differences.

These results do not identify a universal winner. Instead, they show
that the evaluated configurations preserved different properties of the
real data. The Pythia models performed more strongly on predictive
transfer and pairwise linear relationships. HealthGAN performed more
strongly on marginal similarity and Cox confidence-interval overlap.

A possible explanation for the Pythia results is that autoregressive,
class-conditioned generation retained conditional relationships that
were useful for downstream prediction. Because when generating each feature,
model uses both the class label and previously generated features.

HealthGAN's strong marginal  similarity may result from its GAN training
process, which encourages synthetic records to resemble the real data. In GAN training, 
differences in individual feature
distributions can help the discriminator identify synthetic records,
thereby encouraging the generator to reproduce the marginal
distributions of the real data. These explanations are plausible 
but were not tested directly. The
models also differed in architecture, scale, pretraining, representation,
optimisation, and computational cost. The observed differences therefore
cannot be attributed to the generator family alone.

The difference between F1 and AUC-ROC may result from the selected
classification threshold rather than better overall separation of the
two classes. In addition, each configuration was evaluated using a single recorded
run. The stability of the observed model ordering across random seeds was
therefore not tested.

Future studies should repeat the comparison across several clinical
datasets and random seeds. Training data, tuning effort, and
computational budgets should be controlled as well.
Scale-controlled LLM experiments would also help determine whether the
predictive pattern is associated with model size, as suggested by
previous work~\cite{Miletic2024}.

\subsection{RQ3: Observed Changes Under the Private Configurations}

The most informative RQ3 finding comes from the more closely matched
Pythia-1B comparison. The private configuration had lower predictive
performance than the non-private configuration. However, Cox
confidence-interval overlap remained unchanged, and the changes in
statistical similarity were smaller than the changes in predictive
performance. The unchanged Cox overlap is clinically relevant because it
shows that the same level of effect-interval agreement was observed in
both Pythia-1B configurations, even though their predictive results
differed. The reason for that could be that predictive performance 
requires multivariable relationships to be preserved. Statisical 
similarity don't necessarily gurantees that. Also for Cox
confidence-interval doesn't capture difference in interval width or 
the extent of overlap. 

For Pythia-1B-DP, the attribute-inference attack performed no better than
the baseline, and outlier linkability decreased to zero. These results
suggest lower disclosure risk under the selected metrics. However, other
privacy risks and the stability of these results across repeated runs were
not evaluated.

One possible explanation is that private fine-tuning was limited to the
low-rank adapter parameters of the pretrained Pythia-1B model. This may
have allowed the model to retain useful pretrained information while the
adapter updates were clipped and noised. The result is consistent with
previous evidence that larger pretrained language models can perform
better under private fine-tuning~\cite{Li2022LargeDP}. However, the
Pythia-70M and Pythia-1B private runs differed in dataset size, batch
construction, learning rate, and calibrated noise multiplier. Their
different training behaviour cannot therefore be attributed to model
scale alone.

Pythia-1B-DP also retained substantially more measured utility than
PATE-GAN. This cross-family comparison provides weaker evidence because
HealthGAN and PATE-GAN differ in architecture, training objective, and
privacy procedure. PATE-GAN's poor results may reflect its particular
implementation, the dataset, its optimisation settings, or a combination
of these factors. The result should not be interpreted as evidence that
PATE or differentially private GANs generally fail.

A further limitation concerns the scope of the Pythia privacy accounting.
The reported privacy cost applies to the DP-SGD parameter updates. The
loss-based checkpoint comparison was not included in that accounting.
The reported value should therefore not be presented as a complete
end-to-end privacy guarantee for every operation in the training
pipeline.

Future work should compare private and non-private versions of the same
GAN architecture. The analysis should also repeat the Pythia comparison using matched
batch-size and optimisation settings. It should evaluate multiple model
scales and privacy budgets and report uncertainty across repeated runs.
Stronger attacks and privacy-accounted model-selection procedures would
provide a more complete assessment of the relationship between formal
accounting and empirical disclosure risk.

\section{Strengths and Weaknesses of the Two Generator Families}
\label{sec:discussion:gan-vs-llm}

The results show that the two model families have different but
complementary strengths. The strength of the
GAN family is distributional where HealthGAN matches the real marginals more
tightly than any other generator and, once trained, samples cheaply. Its
weaknesses are equally visible here. It enforces no schema-level
constraints, and the only differentially private GAN evaluated does not
produce usable data at the tested budget. However, this result reflects both 
the differential-privacy mechanism
and the change in architecture. It therefore cannot be attributed to
differential privacy alone(\S\ref{sec:discussion:limit:gan-toggle-asymmetry}). 
The strength of the LLM family
is its native handling of mixed-type tabular data through a row-as-text
representation. The same model ingests categorical, ordinal, and continuous
fields without per-column encoders, learns their conditional structure, and
through schema-aware parsing, emits syntactically valid rows. This
representational flexibility underlies the predictive advantage of the LLM
family. Its graceful degradation under differential privacy has a distinct
cause. DP-SGD perturbs a small low-rank adapter at the $1$B
scale, rather than following from the text representation itself. Its
costs are a generation-time parsing-and-retry overhead. It also has substantially longer
training and inference, and a known problemn of word-for-word memorisation
of training text~\cite{Carlini2021Extracting}. Notably, however, the
membership-inference results reported here show no measurable memorisation by
any generator at this scale (\S\ref{sec:results:rq2:privacy:mia}).

\paragraph{Sensitivity to Categorical Granularity:}
The age-binning ablation
(Section~\ref{sec:results:rq2:utility:age-binning}) shows that upstream
categorical-cardinality choices can affect generators differently. Although
the comparison is limited to the two GAN baselines and Pythia-70M, the only
models for which matched three-bin and ten-bin runs exist. Restoring the
native ten-decade \emph{age} encoding raised the predictive utility
of Pythia-70M by $\approx 0.065$ and its F1 by $\approx 0.106$ which is the
largest AUC gain among the compared generators. HealthGAN responded on a
similar scale but in a mixed direction. It gains on F1 ($+0.118$) while
losing slightly on AUC ($-0.029$). PATE-GAN remained weak under
both encodings. No single family is therefore uniformly more sensitive on
this ablation. The practical implication for healthcare data stewards is
that preprocessing decisions
can shift recoverable downstream utility for some generators.

\section{Practical Recommendations for Healthcare Data Stewards}
\label{sec:discussion:recommendations}

The results translate into a small set of evidence-based recommendations
for a data steward choosing a generator for tabular medical data.

\paragraph{When a formal privacy guarantee is required:}
The differentially private LLM (Pythia-1B-DP) is the recommended choice
rather than the differentially private GAN as implemented here. At
$\varepsilon = 5.0$ the DP-LLM retained usable utility (TSTR AUC $0.583$)
and reduced attribute disclosure, whereas the PATE-GAN implementation
evaluated in this thesis collapsed to near chance on every utility metric.
These findings challenge the common use of DP-GAN as the default
approach. However, two caveats apply. First, the DP-LLM must have
sufficient model capacity, as DP-SGD training failed at the $70$M scale.
Second, the evidence for GANs is limited to the PATE-GAN configuration
evaluated in this study (\S\ref{sec:discussion:limit:gan-toggle-asymmetry}).

\paragraph{When maximum utility is required without a formal guarantee:}
Both models are reasonable choices, depending on the main objective.
Pythia-1B provides the highest predictive utility and the strongest
fidelity in joint feature relationships. In contrast, HealthGAN provides
the best marginal fidelity and is the only non-private configuration to
meet or exceed the NNDR reference value of one. It also requires only a
fraction of the computational cost of the LLM.
The choice is the genuine trade-off identified in RQ2 rather than a clear
winner.

\paragraph{When schema handling or operational simplicity dominates:}
The LLM generators ingest mixed-type data without per-column encoders and
emit schema-valid rows directly, which simplifies the pipeline at the cost
of generation time. The GAN requires an explicit encoding and decoding stage
but is far cheaper to run.

\paragraph{Outlier protection:}
No generator concentrated synthetic mass around the vulnerable tail of the
real distribution (outlier linkability $\leq 0.6\%$ for all), so outlier
disclosure was not a discriminating factor in this study. A steward with a
specific outlier-protection requirement should nonetheless verify it
directly, since this property is dataset-dependent.

\section{Toward Standardisation of Synthetic-Data Evaluation}
\label{sec:discussion:standardisation}

The evaluation framework assembled for this thesis, comprising two utility
branches (statistical similarity and predictive performance) and a privacy
pillar reported through twelve complementary single-number metrics. It is
offered as a step toward the reusable evaluation standard called for by
\textcite{hernandez2022synthetic}. General-purpose frameworks such as SynthEval
pursue this goal for arbitrary tabular data~\cite{lautrup2025syntheval}. 
This study applies the same two-part framework to a clinical benchmark
and uses it for a controlled comparison of GAN- and LLM-based approaches. 
Its central principle also applies here: no single metric is sufficient,
and the metrics must be considered together. The PATE-GAN results show
that evaluating fewer dimensions may favour generators that perform well
on one metric while producing data of poor overall quality. What would not
transfer unchanged is the domain-specific content of the suite. The
privacy thresholds (an NNDR of unity, the outlier radius) are calibrated to
the geometry of the real data and must be recomputed per dataset. 
The predictive-utility evaluation requires a predefined outcome
variable, which is used to train and assess the downstream prediction
models. In this study, the target variable is binary readmission. 
The framework should be understood as a reusable \emph{structure} that
jointly evaluates two utility dimensions and privacy. In this study, it
is implemented using metrics suitable for clinical tabular data with a
binary outcome.

\section{Limitations}
\label{sec:discussion:limitations}

\subsection{Single-Dataset Scope}
\label{sec:discussion:limit:dataset}

The entire study is conducted on a single dataset, the Diabetes 130-US
Hospitals dataset. It is a deliberate choice of a large
($> 70{,}000$-record), mixed-type clinical table with binarized
outcome. So it makes it a representative medical tabular testbed, but it is
not universal. The planned evaluation using MIMIC-III was reduced in scope.
Consequently, the findings, particularly the absolute utility values and
model-family rankings, may not generalise to datasets with substantially
different dimensionality, sparsity, or outcome structures.

\subsection{Upper Bound on Language-Model Scale}
\label{sec:discussion:limit:pythia-size}

The LLM-side evaluation is capped at the Pythia-1B scale. 
Pythia-1B was included as a larger non-private and private comparison
point than Pythia-70M. This increase in scale was motivated by the
failure of DP-SGD training at the $70$M scale, as reported in
\S\ref{sec:results:rq3:pythia-70m-failure}. However, larger decoder-only
models with several billion parameters, such as Llama- or Qwen-based
models, were not evaluated. The findings on
the LLM side should therefore be read as characterising the small-to-mid LLM
regime. Model size has been correlated with higher gains, so the
absolute utility figures here are likely a lower bound on what comparable
LLM-based synthesisers can attain at greater scale\cite{Miletic2024} .

\subsection{Architectural Asymmetry of the Cross-Family DP Comparison}
\label{sec:discussion:limit:gan-toggle-asymmetry}

The GAN-family DP comparison reported in
\S\ref{sec:results:rq3:gan-toggle} combines architectural and mechanism
differences. HealthGAN and PATE-GAN are not the same architecture with
differential privacy toggled on or off, but distinct generative designs.
The LLM-family comparison
(\S\ref{sec:results:rq3:llm-toggle}) is, on the contrary, an
architecture-controlled toggle. Pythia-1B and Pythia-1B-DP share the same
decoder-only backbone, the same Tabula serialisation, and the same LoRA
fine-tuning configuration. Although their Fine tuning configuration 
has some differences, but that is due to accommodating DP-SGD in the Pipeline.
So, the observed difference between them is
the DP-SGD training procedure. This asymmetry was difficult to avoid because PATE-GAN has no
standardised non-private counterpart, while applying DP-SGD to WGAN-GP
presents technical challenges, as discussed in
\S\ref{sec:methods:healthgan:nodp}. However, this limitation should be considered when comparing the utility
cost of differential privacy across the two model families. Part of the
GAN-side utility difference reported in
\S\ref{sec:results:rq3:gan-toggle} may result from the architecture change
rather than from differential privacy alone. A further caveat concerns the PATE-GAN
implementation itself. The original PATE-GAN study reported useful 
synthetic-data utility under
stricter privacy budgets than the budget used here~\cite{pategan}.
However, the configuration evaluated in this thesis collapsed on the
current benchmark. This configuration used linear logistic-regression
teachers that were refitted at each training round and a weight-clipped
Wasserstein student
(\S\ref{sec:methods:pategan:teacher}). The observed collapse should therefore be read
as a property of this implementation, dataset, and budget, not as evidence
that differentially private GANs fail in general.

\subsection{Dual Role of PATE-GAN in the Landscape Comparison}
\label{sec:discussion:limit:pategan-dual-role}

PATE-GAN appears alongside non-private models in the RQ2 landscape view
(\S\ref{sec:results:rq2:landscape}) even though it incorporates differential
privacy by construction. This is a reporting choice rather than a flaw.
PATE-GAN is the de-facto DP-GAN baseline in the literature~\cite{pategan},
and excluding it from the landscape figure would make it difficult to determine where the GAN
family lands once formal DP is required. Its position on the privacy-utility
plane should be read accordingly. It is the DP-enforced operating point of
the GAN family, whereas the non-private generators (HealthGAN, Pythia-70M,
and Pythia-1B) sit on the plane without any formal guarantee.

\subsection{Single Privacy Budget ($\varepsilon = 5.0$)}
\label{sec:discussion:limit:epsilon}

Both differentially private generators are evaluated at a single privacy
budget, $\varepsilon = 5.0$ with $\delta = 10^{-5}$. This fixes the
comparison at one point on the privacy-utility curve and cannot reveal its
slope. A budget sweep (for example $\varepsilon \in \{1, 3, 5, 10\}$) would
show whether the graceful degradation of the LLM family persists at tighter
budgets and whether the GAN family becomes competitive at looser ones. This
is left to future work.

\subsection{Computational Cost}
\label{sec:discussion:limit:compute}

The computational costs of the two paradigms differ by orders of magnitude
and were not equalised. The dp variant of LLM dominates here. In this implementation, Opacus computes gradients for individual
training examples. This requires more memory than standard mini-batch
SGD and therefore required the effective-batch-size strategy described
in \S\ref{sec:results:rq3:llm-toggle}. The GAN baseline, by contrast, trains
and samples cheaply. Comparisons in this thesis are made at a fixed privacy
budget rather than at fixed compute. The utility rankings should be read
with that asymmetry in mind.

\section{Threats to Validity}
\label{sec:discussion:validity}
Four types of validity threat are considered. \emph{Internal validity}
is limited by configuration differences. It applies particularly in the GAN
comparison where architecture and privacy change together.
\emph{External validity} is limited by the use of one dataset, one
privacy budget, and one decoder-only model family. The rankings may not
generalise to other settings. \emph{Construct validity} is limited
because no single metric fully captures privacy or utility. For example,
PATE-GAN achieved favourable distance-based privacy scores despite poor
utility. \emph{Conclusion validity} is limited by the use of one data
split and one random seed, resulting in point estimates without
uncertainty intervals. The attribute-inference results should therefore
be interpreted as suggestive rather than definitive.

\section{Ethical Considerations}
\label{sec:discussion:ethics}

The motivation for synthetic medical data is itself ethical, namely to
enable analysis without exposing patient records. But the practice carries
its own obligations. Even differentially private synthetic data is generated
from real patient records collected under specific consent and data-use
agreements. The synthetic output also does not inherit blanket permission to
be shared without regard to those terms. A particular hazard is false
confidence. A formal $\varepsilon$ guarantee bounds one specific privacy
notion and does not certify that every disclosure channel is closed. Also, as
the PATE-GAN case shows, strong-looking privacy numbers can instead reflect
unusable output. Releasing a trained generator rather than the data shifts
but does not remove this risk, since the model encodes information about its
training set and remains susceptible to extraction attacks in
principle~\cite{Carlini2021Extracting}. Finally, synthetic data that
faithfully preserves the statistical structure of a clinical population also
reproduces the biases embedded in that population. So downstream models
trained on it warrant the same fairness scrutiny as models trained on the
real data.

\cleardoublepage
